\hnsection{指数映射与豪斯多夫公式}

关于指数映射的最初研究由 E. Study 和 F. Engel 完成；Engel {[}9 b{]} 注意到，对于 \(\mathbf{SL}(2, \mathbf{C})\)，指数映射不是满射（例如 \(\begin{pmatrix} -1 & a \\ 0 & -1 \end{pmatrix}\) 在 \(a \neq 0\) 时不是指数元），但对于 \(\mathbf{GL}(n, \mathbf{C})\) 是满射，因而对于 \(\mathbf{PGL}(n, \mathbf{C})\) 也是满射（后一个性质在 \(n = 2\) 时 Study 已经注意到）；由此，\(\mathbf{SL}(2, \mathbf{C})\) 和 \(\mathbf{PGL}(2, \mathbf{C})\) 给出了两个局部同构群的例子，然而从全局观点看它们却非常不同。Engel 还证明，其他经典群加上伸缩变换后指数映射也是满射；Maurer、Study 等人继续开展了这项工作，但没有产生实质性的新结果。

1899 年，H. Poincaré ({[}14{]}, vol.~3, pp.~169--172 and 173--212) 从另一个角度开始研究指数映射。他的论文似乎经过仓促编辑，因为他在数处断言连通群的每个元素都是指数元，而在别处却给出了相反的例子。他的结果主要涉及伴随表示：他证明，这样的群 G 的一个半单元素可能是李代数 L(G) 中无穷多个元素的指数，而一个非半单元素则可能根本不是指数元。如果 ad(X) 没有作为 2πi 的非零倍数的特征值，那么 exp 在 X 处是 étale 的。他还证明，若 U 和 V 在 L(G) 中描出回路，W 由连续性定义使得 \(e^{U} \cdot e^{V} = e^{W}\) ，也可能无法回到 W 的原始值。他使用了一个残数公式，其本质上相当于

\[
\Phi (\mathrm{adX}) = \frac {1}{2 \pi i} \int \frac {\Phi (\xi) d \xi}{\xi - \mathrm{adX}}
\]

其中 \(\operatorname{ad}(\mathbf{X})\) 是一个非零特征值重数均为 1 的半单元素，\(\Phi\) 是收敛半径足够大的整级数，积分沿包围 \(\operatorname{ad} X\) 的特征值的闭合曲线取；他还研究了当 X 趋近于具有重特征值的变换时所发生的情况。

关于在公式 \(e^{U}.e^{V}=e^{W}\) 中把 W 表示为 U 和 V 的函数的研究，在 Poincaré 的工作之前不久就已成为 Campbell 的两篇论文的主题 {[}13{]}。正如 Baker 不久后所写：“……李理论以显然的方式表明，乘积 \(e^{U}.e^{V}\) 的形式是 \(e^{W}\)，其中 W 是 U 和 V 的交错式级数……”。后来关于这一问题的工作旨在使这一断言精确化，并给出 W 的显式公式（或构造方法）（“豪斯多夫公式”）。继 Campbell 和 Poincaré 之后，Pascal、Baker {[}15{]} 和 Hausdorff {[}16{]} 又回到这个问题；他们各自都认为前人的证明并不令人信服；主要困难在于“交错式”究竟是什么意思：它们是所讨论李代数的特殊元素，还是普适的“符号”表达式？Campbell、Poincaré 和 Baker 都没有清楚地说明这一点。另一方面，Hausdorff 的论文极其精确；他首先研究有限个未定元的（非交换）结合形式幂级数代数，并把 U、V、W 看作这个代数的元素。他通过类似其前人的微分方程论证证明了 W 的存在。当未定元替换为有限维李代数的元素时，他又用同一论证证明了级数的收敛性。正如 Baker 所指出的（Poincaré 也独立地指出了这一点），这个结果可用于证明 Lie 的第三定理；他澄清了李群与李代数之间的对应关系，例如在换位子群方面。

1947 年，Dynkin {[}39{]} 再次研究了这个问题，并通过一开始就考察赋范李代数（无论有限维与否，定义在 R、C 或超度量域上），得到了豪斯多夫公式的显式系数。†

